\chapter{The Structural DAG of the URF Framework}

\section{Overview}

The Unified Rigidity Framework (URF) can be organized as a directed acyclic graph (DAG) of structural dependencies.  
Each node represents a theorem, lemma, invariant, or structural principle.  
An edge $A \to B$ indicates that result $B$ depends logically on $A$.

\section{Nodes}

We organize the core nodes as follows:

\begin{itemize}
\item \textbf{Locality Axiom} — Refinement operators are $R$-local.
\item \textbf{Capacity Bound} — Each activation extracts at most $c$ bits per vertex.
\item \textbf{$R$-Power Graph} — $G_n^{(2R)}$ encoding locality constraints.
\item \textbf{$R$-Packing Number} — $\alpha_R(G_n)$.
\item \textbf{$R$-Chromatic Number} — $\chi_R(G_n)$.
\item \textbf{Packing Bound} — $\alpha_R(G_n) \ge n/B_\Delta(2R)$.
\item \textbf{Covering Bound} — $T \ge \chi_R(G_n)$.
\item \textbf{Information Bound} — $T \ge H(X_n)/(c\,\alpha_R(G_n))$.
\item \textbf{Unified Depth Theorem} — 
\[
T \ge \max\!\left\{
\frac{H(X_n)}{c\,\alpha_R(G_n)},
\;
\chi_R(G_n)
\right\}.
\]
\end{itemize}

\section{Dependency Structure}

The logical flow is:

\begin{enumerate}
\item Locality Axiom $\Rightarrow$ $R$-Power Graph.
\item $R$-Power Graph $\Rightarrow$ definitions of $\alpha_R$ and $\chi_R$.
\item $\alpha_R$ $\Rightarrow$ Information Bound.
\item $\chi_R$ $\Rightarrow$ Covering Bound.
\item Information Bound + Covering Bound $\Rightarrow$ Unified Depth Theorem.
\end{enumerate}

\section{DAG Representation}

\begin{figure}[h]
\centering
\begin{tabular}{cc}
Locality Axiom & Capacity Bound \\
$\downarrow$ & $\downarrow$ \\
$R$-Power Graph & Information Bound \\
$\downarrow$ & \\
$\alpha_R,\chi_R$ & \\
$\downarrow$ & \\
Covering Bound & \\
$\downarrow$ & \\
Unified Depth Theorem & \\
\end{tabular}
\caption{Structural dependency DAG for refinement depth lower bounds.}
\label{fig:urf-dag}
\end{figure}

\section{Structural Interpretation}

The DAG shows that depth lower bounds arise from two independent mechanisms:

\begin{itemize}
\item \emph{Packing (information) constraint} — limits parallel entropy extraction.
\item \emph{Covering (chromatic) constraint} — limits simultaneous vertex activation.
\end{itemize}

The Unified Depth Theorem is the minimal upper envelope of these two constraints.

\section{Acyclicity}

No result in the diagram depends on the Unified Depth Theorem itself.  
All arrows point strictly forward from axioms to consequences, ensuring the structure is acyclic.

\section{Implication}

Any improvement to depth bounds must refine either:

\begin{itemize}
\item the packing parameter $\alpha_R(G_n)$,
\item the chromatic parameter $\chi_R(G_n)$, or
\item the capacity constant $c$.
\end{itemize}

Thus the DAG isolates the exact structural levers controlling refinement depth.
